\scanpage{214}
\mspage{1 - 49 -}\hfill n\textsuperscript{o} 413

Soient $\mathrm G$ et $\mathrm H$ deux noyaux de groupe de classe $\mathrm C^{\mathrm r}$. Un \uline{morphisme} $\psi$ de $\mathrm G$ dans $\mathrm H$ est une application de classe $\mathrm C^{\mathrm r}$ de $\mathrm G$ dans $\mathrm H$ satisfaisant aux conditions suivantes :

(a) $\psi(\mathrm e_{\mathrm G})=\mathrm e_{\mathrm H}\qquad\psi(\unsure{\vartheta|\theta}_{\mathrm G}(\mathrm x))=\vartheta_{\mathrm H}(\psi(\mathrm x))$ pour tout $\mathrm x\in\mathrm G$ ;

(b) si $(\mathrm x,\mathrm y)\in\Omega(\mathrm G)$, alors $(\psi(\mathrm x),\psi(\mathrm y))\in\Omega(\mathrm H)$ et on a :
\[
\mathrm m_{\mathrm H}(\psi(\mathrm x),\psi(\mathrm y))=\psi(\mathrm m_{\mathrm G}(\mathrm x,\mathrm y))
\]
(autrement dit $\psi$ est une représentation pour les lois de composition non partout définies sur $\mathrm G$ et $\mathrm H$). Nous laisserons au lecteur le soin d'étendre au cas des noyaux de groupe les propositions 3 et 3 bis.

Soit $\mathrm G$ un noyau de groupe. Il existe un voisinage ouvert $\mathrm V$ de $\mathrm e$ tel que les conditions de (iii) soient réalisées pour $\mathrm x,\mathrm y,\mathrm z\in\mathrm V$. En particulier le \uline{produit de convolution} associé à la restriction de $\mathrm m$ à $\mathrm V\times\mathrm V$ définit une loi de composition sur $\mathrm T_{\mathrm e}^{(\mathrm r)}(\mathrm G)$, qui est bilinéaire et associative (bien que non partout définie si $\mathrm r<\infty$). Nous désignerons encore par $\mathrm U(\mathrm G)$ l'espace $\mathrm T_{\mathrm e}^{(\mathrm r)}(\mathrm G)$ muni de cette multiplication. Nous laisserons au lecteur le soin d'étendre les autres notations et résultats des n\textsuperscript{os} 12 à 16. Signalons par exemple que si $\mathrm f$ est un morphisme de $\mathrm G$ dans un noyau de groupe $\mathrm H$, alors $\mathrm f_{*}$ est une \uline{représentation} de $\mathrm U(\mathrm G)$ dans $\mathrm U(\mathrm H)$.

De plus, l'application tangente $\mathrm m^{\mathrm T}$ à $\mathrm m$ est une application de classe $\mathrm C^{\mathrm r-1}$ de $\mathrm T(\quad)$ dans $\mathrm T(\mathrm G)$. On peut identifier $\mathrm T(\quad)$ à un ouvert de $\mathrm T(\mathrm G)\quad\mathrm T(\mathrm G)$ et on vérifie aisément que $\mathrm T(\mathrm G)$ muni du point marqué $0_{\mathrm e}$, de l'application $\mathrm m^{\mathrm T}$ de $\mathrm T(\quad)$ dans $\mathrm T(\mathrm G)$ et de l'application $\vartheta^{\mathrm T}$ de $\mathrm T(\mathrm G)$ dans $\mathrm T(\mathrm G)$ est un \uline{noyau de groupe} de classe $\mathrm C^{\mathrm r-1}$ ; la projection canonique est un morphisme de classe $\mathrm C^{\mathrm r-1}$ de $\mathrm T(\mathrm G)$ sur $\mathrm G$. Nous laisserons au lecteur le soin d'étendre au cas des noyaux de groupe les propositions 11 et 12 ainsi que leurs corollaires.
% archive sheet number: 213

